\chapter{Solitons and Collective Quantization}
\label{ch:solitons-collective-quantization}
\ConventionsLorentzian

The finite-energy solitons of
Chapter~\ref{ch:vol2-classical-yang-mills-matter} are classical saddles in
topological sectors.  Quantization adds three pieces which are invisible in
the static solution itself: collective coordinates for exact zero modes,
determinants on the orthogonal fluctuation space, and the definition of the
physical observable or Hilbert-sector coordinate.  This chapter starts the
dedicated soliton quantization lane by keeping those three ingredients in one
calculation.

The guiding example is deliberately elementary.  A kink on the line has a
one-dimensional moduli space of centers.  The center is easy; the mass of the
soliton particle is not fixed by that moduli space.  It is fixed by removing
the translation zero mode, computing the fluctuation determinant in the
chosen renormalization convention, and then identifying the isolated pole or
sector energy.  This is the same physics lesson as for instantons, now in a
Lorentzian sector with a soliton particle instead of a Euclidean tunneling
event.

\section{Kink Sector and the Translation Collective Coordinate}
\label{sec:kink-sector-translation-collective-coordinate}

Consider the dimensionless scalar theory in \(1+1\) dimensions with static
energy
\begin{equation}
  E[\phi]
  =
  \int_{\mathbb R}\dd x\,
  \left[
    \frac12(\partial_x\phi)^2
    +\frac12(1-\phi^2)^2
  \right].
  \label{eq:phi4-kink-static-energy}
\end{equation}
The two vacua are \(\phi=\pm1\).  Finite-energy configurations are separated
by the boundary coordinate
\[
  Q_{\rm kink}
  =
  \frac12\bigl(\phi(+\infty)-\phi(-\infty)\bigr)\in\mathbb Z
\]
for smooth fields approaching the vacua at infinity.  The charge-one kink is
\begin{equation}
  \phi_K(x;X)=\tanh(x-X),
  \qquad
  \partial_x\phi_K=\operatorname{sech}^2(x-X).
  \label{eq:phi4-kink-profile}
\end{equation}

\begin{proposition}[Kink center as a collective coordinate]
\label{prop:kink-center-collective-coordinate}
For the kink \eqref{eq:phi4-kink-profile}, the translation zero-mode norm is
equal to the classical kink mass:
\begin{equation}
  G_{XX}
  =
  \int_{\mathbb R}\dd x\,(\partial_X\phi_K)^2
  =
  \int_{\mathbb R}\dd x\,(\partial_x\phi_K)^2
  =
  M_{\rm cl}
  =
  \frac43 .
  \label{eq:kink-center-collective-metric}
\end{equation}
If
\[
  \phi(t,x)=\phi_K(x-X(t))+\eta(t,x),
  \qquad
  \int_{\mathbb R}\dd x\,\eta(t,x)\,\partial_x\phi_K(x-X(t))=0 ,
\]
then the leading small-velocity kinetic term is
\begin{equation}
  \frac12\int_{\mathbb R}\dd x\,\dot\phi^2
  =
  \frac12M_{\rm cl}\dot X^2
  +
  \frac12\int_{\mathbb R}\dd x\,\dot\eta^2
  +\hbox{higher terms in }\eta,\dot X .
  \label{eq:kink-collective-coordinate-kinetic-split}
\end{equation}
Thus the center coordinate has the inertial mass \(M_{\rm cl}\), and the
Gaussian determinant for normal modes must be taken on the subspace
orthogonal to \(\partial_x\phi_K\).
\end{proposition}

\begin{proof}
The kink satisfies the first-order equation
\[
  \partial_x\phi_K=1-\phi_K^2.
\]
Therefore the static energy density is
\[
  \frac12(\partial_x\phi_K)^2+\frac12(1-\phi_K^2)^2
  =
  (\partial_x\phi_K)^2
  =
  \operatorname{sech}^4(x-X).
\]
The integral
\[
  \int_{-\infty}^{\infty}\operatorname{sech}^4 u\,\dd u
  =
  \int_{-1}^{1}(1-y^2)\,\dd y
  =
  \frac43,
  \qquad y=\tanh u ,
\]
gives both \(M_{\rm cl}\) and the zero-mode norm.  Since
\(\partial_X\phi_K=-\partial_x\phi_K\), the term linear in
\(\dot X\dot\eta\) is proportional to
\(\int\dot\eta\,\partial_x\phi_K\), which is zero after differentiating the
orthogonality condition at leading order.  This gives
\eqref{eq:kink-collective-coordinate-kinetic-split}.
\end{proof}

The formula \(\sqrt{G_{XX}}\,\dd X\) is the collective-coordinate measure for
the center.  It is not the quantum mass shift.  The latter comes from the
determinant of the normal fluctuation operator after the zero mode has been
removed.

\section{Fluctuation Determinants and the Soliton Mass}
\label{sec:soliton-fluctuation-determinants-mass}

The sine-Gordon model gives the cleanest benchmark because the exact
soliton-breather spectrum is independently known in
Chapter~\ref{ch:integrable-sine-gordon-thirring-affine-toda}.  In the
normalization of Section~\ref{sec:sine-gordon-dhn-soliton-mass}, the kink
has classical mass \(8m/\beta^2\).  Its normal fluctuation operator has a
translation zero mode and continuum phase shift
\begin{equation}
  \delta(k)=2\arctan(m/k),
  \qquad
  \frac{\dd\delta}{\dd k}
  =
  -\frac{2m}{k^2+m^2}.
  \label{eq:soliton-chapter-sg-phase-shift}
\end{equation}

\begin{proposition}[DHN mass shift as a determinant statement]
\label{prop:soliton-chapter-dhn-mass-shift}
In the Dashen--Hasslacher--Neveu/Faddeev--Korepin normal-ordering convention,
the sine-Gordon one-loop soliton mass shift is
\begin{equation}
  \Delta M_{\rm SG}^{(1)}
  =
  \frac12
  \int_{-\Lambda}^{\Lambda}
  \frac{\dd k}{2\pi}\,
  \sqrt{k^2+m^2}\,
  \frac{\dd\delta}{\dd k}
  +
  \Delta M_{\rm ct}(\Lambda)
  =
  -\frac{m}{\pi},
  \label{eq:soliton-chapter-dhn-mass-shift}
\end{equation}
where
\[
  \Delta M_{\rm ct}(\Lambda)
  =
  \frac{m}{\pi}\operatorname{arsinh}(\Lambda/m)-\frac{m}{\pi}.
\]
Consequently
\begin{equation}
  M_s
  =
  \frac{8m}{\beta^2}-\frac{m}{\pi}+O(\beta^2m).
  \label{eq:soliton-chapter-sg-physical-mass}
\end{equation}
\end{proposition}

\begin{proof}
Substituting \eqref{eq:soliton-chapter-sg-phase-shift} into the continuum
determinant comparison gives
\[
  \frac12
  \int_{-\Lambda}^{\Lambda}
  \frac{\dd k}{2\pi}
  \sqrt{k^2+m^2}
  \left(-\frac{2m}{k^2+m^2}\right)
  =
  -\frac{m}{\pi}\operatorname{arsinh}(\Lambda/m).
\]
The counterterm displayed above is the same tadpole subtraction which defines
the normal-ordered cosine.  The logarithmic cutoff term cancels, leaving the
finite value \(-m/\pi\).  The translation zero mode is absent from the
determinant because it was converted into the center coordinate in
Proposition~\ref{prop:kink-center-collective-coordinate}.
\end{proof}

The finite term in \eqref{eq:soliton-chapter-dhn-mass-shift} is not a
decoration of the classical kink.  It is the difference between the mass
parameter entering the exact soliton scattering theory and the mass obtained
by quantizing only the moduli-space center.  The exact lightest-breather
check in Proposition~\ref{prop:sine-gordon-dhn-breather-mass-consistency}
uses precisely this finite determinant term.

\section{Fermion Zero Modes and Fractional Sector Charge}
\label{sec:soliton-fermion-zero-modes-fractional-charge}

Soliton quantization can also change the charge carried by a sector.  In the
Jackiw--Rebbi kink model, the one-particle Dirac Hamiltonian may be written
as
\begin{equation}
  H_{\rm JR}
  =
  \begin{pmatrix}
    0 & -\partial_x+M(x)\\
    \partial_x+M(x) & 0
  \end{pmatrix},
  \qquad
  M(x)=m\tanh(mx).
  \label{eq:soliton-chapter-jackiw-rebbi-kink-hamiltonian}
\end{equation}
The kink mass profile changes sign between the two vacua.  The normalized
zero mode is
\begin{equation}
  \psi_0(x)
  =
  \sqrt{\frac m2}
  \begin{pmatrix}
    \operatorname{sech}(mx)\\
    0
  \end{pmatrix},
  \qquad
  H_{\rm JR}\psi_0=0 .
  \label{eq:soliton-chapter-jackiw-rebbi-kink-zero-mode}
\end{equation}
All nonzero modes are paired by charge conjugation.  The soliton-sector
fermion number therefore differs by one unit depending on whether this single
zero mode is empty or filled:
\begin{equation}
  F_{\rm kink}
  =
  N_0-\frac12,
  \qquad
  N_0\in\{0,1\}.
  \label{eq:soliton-chapter-jackiw-rebbi-kink-half-charge}
\end{equation}

This is a physical sector statement, not a statement about the location of
the kink center.  The center coordinate organizes translations of the soliton
particle.  The fermion zero mode organizes an internal two-state sector whose
charges are \(-1/2\) and \(+1/2\).  A semiclassical soliton calculation is
complete only after both sorts of zero modes have been separated and the
nonzero fluctuation determinant is paired with the observable being measured.
